\documentclass[aspectratio=169,11pt]{beamer}
\usepackage[indonesian]{babel}
\usepackage{booktabs,tikz,listings}
\usetheme{metropolis}
\usebeamercolor{normal text}
\colorlet{codeink}{normal text.fg}
\usebeamercolor{alerted text}
\colorlet{codeaccent}{alerted text.fg!60!codeink}
\lstset{language={[LaTeX]TeX},basicstyle=\ttfamily\small\color{codeink},
  texcsstyle=*\bfseries\color{codeaccent},moretexcs={draw,fill,centering,caption,label,ref,cite,approx,mathrm},
  commentstyle=\itshape\color{codeink!75},columns=fullflexible,keepspaces=true,
  showstringspaces=false,backgroundcolor=\color{codeink!6},frame=l,
  rulecolor=\color{codeaccent},xleftmargin=5pt,aboveskip=4pt,belowskip=4pt}
\lstdefinestyle{terminal}{language=bash,morekeywords={latexmk},keywordstyle=\bfseries\color{codeaccent}}
\newcommand{\tugas}[1]{\par\medskip\alert{\textbf{Coba sendiri}}\quad #1}
\title{Menulis dokumen akademik dengan \LaTeX}
\subtitle{Hari 1 --- Dasar penulisan dan manajemen proyek}
\author{Pelatihan Universitas Terbuka dan Departemen Matematika FMIPA UGM}
\date{Yogyakarta, 19 November 2026}
\begin{document}
\begin{frame}[plain]
\usebeamerfont{title}\inserttitle\par
\usebeamerfont{subtitle}\insertsubtitle\par\medskip
\usebeamertemplate{title separator}\par\medskip
\usebeamerfont{author}\insertauthor\par\smallskip
\usebeamerfont{date}\insertdate
\end{frame}
\begin{frame}[fragile]{Pertanyaan yang akan kita jawab}
Jika panjang tali dilipatduakan, apakah periode ayunan juga menjadi dua kali lebih lama?
\medskip

Hari ini Anda akan menyusun laporan pendek yang memuat dugaan, model matematika, data, gambar alat, dan kesimpulan. Simpan sumbernya. Pada hari kedua, hasil ini menjadi bahan presentasi Anda.
\end{frame}
\begin{frame}{Peta belajar Hari 1}
\begin{enumerate}
\item \textbf{Pengenalan, struktur, dan format dokumen}\\
Membuat sumber yang dapat dikompilasi dan mudah dibaca.
\item \textbf{Elemen akademik}\\
Menulis rumus, tabel, gambar, dan rujukan otomatis.
\item \textbf{Proyek dokumen akademik lengkap}\\
Menggabungkan seluruh elemen menjadi laporan pendek.
\item \textbf{Manajemen proyek dokumen akademik besar}\\
Memisahkan berkas dan mengelola detail secara konsisten.
\end{enumerate}
\end{frame}
\section{Sesi 1. Pengenalan, struktur, dan format dokumen}
\begin{frame}[fragile]{Pengenalan, struktur, dan format dokumen}
\relax{\large Sesi 1 \quad 08.30--09.45 \quad 75 menit}\par\bigskip
10 menit mengenali LaTeX, editor, dan alur kompilasi\\[3pt]
15 menit membuat dokumen, kelas, dan prambule\\[3pt]
35 menit menyusun bagian, paragraf, format, daftar, dan catatan\\[3pt]
15 menit membaca log dan memperbaiki kesalahan dasar
\end{frame}
\begin{frame}{Mengapa menggunakan LaTeX?}
\begin{itemize}
\item Struktur ditulis dalam sumber; tata letak mengikuti kelas dokumen.
\item Rumus, penomoran, rujukan silang, dan daftar pustaka dapat dikelola secara konsisten.
\item Sumber teks mudah dipisahkan, ditinjau, dan digunakan kembali.
\end{itemize}
\medskip
\textbf{Alur kerja}\quad sumber \texttt{.tex} $\rightarrow$ kompilasi
$\rightarrow$ PDF dan log $\rightarrow$ revisi sumber.
\medskip

Overleaf menyediakan editor dan kompilasi di peramban. Secara lokal,
Anda memerlukan editor serta distribusi TeX.
\end{frame}
\begin{frame}[fragile]{Proyek pertama}
Di Overleaf, buat proyek kosong dan buka \texttt{main.tex}. Ketik contoh pada slide berikut, lalu tekan \textit{Recompile}.

\medskip
Jika bekerja secara lokal, simpan \texttt{main.tex} dalam folder proyek. Dengan distribusi TeX dan \texttt{latexmk} terpasang, jalankan di folder itu:
\begin{lstlisting}[style=terminal]
latexmk -pdf main.tex
\end{lstlisting}
Buka \texttt{main.pdf}. Jika kompilasi terhenti, pesan pertama pada log biasanya menunjukkan tempat terbaik untuk mulai memeriksa sumber.
\end{frame}
\begin{frame}[fragile]{Berkas pertama yang dapat dikompilasi}
\begin{lstlisting}
\documentclass[11pt,a4paper]{article}
\usepackage[indonesian]{babel}
\title{Mengukur periode ayunan}
\author{Nama peserta}
\date{}
\begin{document}
\maketitle
Berapa lama satu getaran berlangsung?
\end{document}
\end{lstlisting}
\tugas{Ganti nama penulis dan kompilasi lagi. Temukan perubahan itu di PDF.}
\end{frame}
\begin{frame}[fragile]{Di mana sebuah perubahan ditulis?}
\begin{description}
\item[Kelas] \lstinline|article| menyediakan bagian dan subbagian. Untuk buku yang memerlukan bab, kelas \lstinline|book| menyediakan \lstinline|\chapter|.
\item[Prambule] Bagian sebelum \lstinline|\begin{document}| memuat paket dan identitas naskah.
\item[Isi] Perintah \lstinline|\maketitle| mencetak identitas. Teks dan lingkungan seperti \lstinline|enumerate| ditempatkan di sini.
\end{description}
\tugas{Ubah opsi kelas dari \texttt{11pt} ke \texttt{12pt}. Bandingkan panjang baris pertama.}
\end{frame}
\begin{frame}[fragile]{Kelas, paket, perintah, dan lingkungan}
\begin{lstlisting}
\documentclass{article}        % kelas dokumen
\usepackage{amsmath}          % paket tambahan
% Isi berikut berada setelah \begin{document}
\textbf{Hasil pengamatan}     % perintah + argumen
\begin{itemize}               % awal lingkungan
  \item Periode bertambah.
\end{itemize}                 % akhir lingkungan
\end{lstlisting}
Kelas berasal dari berkas \texttt{.cls}; paket biasanya dari \texttt{.sty}.
Muat paket di prambule. Pasangkan setiap lingkungan pembuka dan penutup.
\end{frame}
\begin{frame}[fragile]{Hierarki bagian dan bab}
\begin{lstlisting}
\section{Pendahuluan}
\subsection{Pertanyaan penelitian}
\subsubsection{Asumsi model}
\section{Metode}
\section{Hasil dan pembahasan}
\section{Kesimpulan}
\end{lstlisting}
\lstinline|article| menggunakan bagian. Kelas \lstinline|report| dan
\lstinline|book| juga menyediakan \lstinline|\chapter|.
Versi berbintang, seperti \lstinline|\section*{Ucapan terima kasih}|,
tidak diberi nomor.
\tugas{Buat empat bagian utama laporan sebelum mengisi paragrafnya.}
\end{frame}
\begin{frame}[fragile]{Paragraf dan penekanan dalam naskah}
\begin{lstlisting}
Kami mengukur waktu untuk \textbf{sepuluh getaran}.

Satu getaran berakhir ketika beban kembali ke titik
awal dengan \emph{arah gerak yang sama}.
\end{lstlisting}
Kami mengukur waktu untuk \textbf{sepuluh getaran}.

\smallskip
Satu getaran berakhir ketika beban kembali ke titik awal dengan \emph{arah gerak yang sama}.

\tugas{Hapus baris kosong, lalu kompilasi. Apa yang berubah?}
\end{frame}
\begin{frame}[fragile]{Baris baru, spasi, dan karakter khusus}
\begin{lstlisting}
Baris sumber pertama
masih menjadi paragraf yang sama.

Ini paragraf baru. Nilai galat adalah 2\%.
Tabel~1 tidak memisahkan kata dan nomor.
Simbol khusus: \& \_ \# \$ \{ \}
\end{lstlisting}
Beberapa spasi biasa dianggap satu spasi. Satu baris kosong memulai
paragraf baru. \lstinline|\\| memutus baris dalam konteks yang sesuai,
misalnya alamat; jangan gunakan sebagai pengganti paragraf.
\medskip

Tanda \lstinline|%| memulai komentar sampai akhir baris sumber.
\end{frame}
\begin{frame}[fragile]{Penekanan dan perataan teks}
\begin{lstlisting}
\textbf{Tebal}, \textit{miring}, \underline{garis bawah}.
\emph{Penekanan mengikuti konteks jenis huruf.}
\begin{center}
  Laporan pengukuran periode ayunan
\end{center}
\begin{flushleft}
  Kelompok A
\end{flushleft}
\end{lstlisting}
Gunakan \lstinline|flushright| untuk rata kanan. Teks biasa pada
\lstinline|article| diratakan kiri dan kanan secara otomatis.
\tugas{Beri penekanan pada istilah penting tanpa memformat seluruh paragraf.}
\end{frame}
\begin{frame}[fragile]{Alat, bahan, dan langkah kerja}
\begin{lstlisting}
\section{Metode}
\begin{itemize}
  \item Tali, beban, dan penyangga
  \item Mistar dan pencatat waktu
\end{itemize}
\begin{enumerate}
  \item Ukur panjang dari titik gantung ke pusat beban.
  \item Catat waktu untuk sepuluh getaran.
\end{enumerate}
\end{lstlisting}
\tugas{Tambahkan satu langkah agar kelompok lain dapat mengulang pengukuran dengan cara yang sama.}
\end{frame}
\begin{frame}[fragile]{Daftar istilah dan catatan tambahan}
\begin{lstlisting}
\begin{description}
  \item[Periode] Waktu untuk satu getaran.
  \item[Panjang] Jarak titik gantung ke pusat beban.
\end{description}
Data ini bersifat ilustratif.\footnote{Bukan hasil eksperimen.}
Periksa pengukuran awal.\marginpar{Ukur ulang}
\end{lstlisting}
\lstinline|\footnote| menambah catatan kaki; \lstinline|\marginpar|
menambah catatan margin pada dokumen artikel.
Catatan margin memerlukan ruang dan dapat bertumpuk.
\tugas{Tambahkan satu catatan kaki yang menjelaskan asal data.}
\end{frame}
\begin{frame}[fragile]{Membaca kesalahan kompilasi}
\begin{description}
\item[Perintah tidak dikenal] Periksa ejaan perintah dan paket yang diperlukan.
\item[Kurung atau lingkungan] Cocokkan \lstinline|{| dengan \lstinline|}|,
serta \lstinline|\begin| dengan \lstinline|\end|.
\item[Mode matematika] Indeks dan pangkat harus berada di dalam rumus.
\item[Urutan perbaikan] Baca kesalahan pertama pada log, periksa juga baris
sebelumnya, ubah satu hal, lalu kompilasi kembali.
\end{description}
\texttt{Overfull} adalah peringatan teks melewati ruang yang tersedia;
berbeda dari kesalahan yang menghentikan kompilasi.
\end{frame}
\begin{frame}[fragile]{Uji kompilasi berpasangan}
Pada salinan berkas, ubah \lstinline|\end{enumerate}| menjadi \lstinline|\end{itemize}|. Prediksi hasilnya, lalu kompilasi.

\medskip
Tunjukkan baris pertama yang dilaporkan bermasalah kepada rekan. Cocokkan nama lingkungan pembuka dan penutup, perbaiki, lalu kompilasi lagi.

\medskip
\textbf{Hasil sesi pertama}\quad PDF menampilkan identitas penulis, pertanyaan penelitian, daftar alat, dan langkah bernomor.
\end{frame}
\section{Sesi 2. Elemen dokumen akademik}
\begin{frame}[fragile]{Elemen dokumen akademik}
\relax{\large Sesi 2 \quad 10.00--11.30 \quad 90 menit}\par\bigskip
25 menit menulis model periode dan persamaan bernomor\\[3pt]
25 menit menyusun tabel dari data ilustratif\\[3pt]
20 menit menggambar alat dan merujuk gambar\\[3pt]
20 menit menyisipkan sitasi dan memeriksa daftar pustaka
\end{frame}
\begin{frame}[fragile]{Panjang tali dan periode ayunan}
\begin{columns}[c]\column{.37\linewidth}
\centering\begin{tikzpicture}[scale=1.1]
\draw[thick] (-1,0)--(1,0);
\draw[dashed] (0,0)--(0,-2.5);
\draw[thick] (0,0)--(-65:2.4) node[midway,right] {$L$};
\fill (-65:2.4) circle (.13);
\draw (0,-.7) arc(-90:-65:.7);
\node at (-78:1.03) {$\theta$};
\end{tikzpicture}
\column{.58\linewidth}
Pada simpangan kecil, beban mengikuti model
\[T\approx2\pi\sqrt{\frac{L}{g}}.\]
$L$ adalah jarak titik gantung ke pusat beban. $T$ adalah waktu satu getaran. Anggap tali ringan dan redaman kecil.
\end{columns}
\medskip
Jika $L$ menjadi $2L$, model memprediksi periode menjadi $\sqrt{2}$ kali periode awal.
\end{frame}
\begin{frame}[fragile]{Rumus di dalam kalimat dan rumus bernomor}
\begin{lstlisting}
Satu periode adalah $T=t/N$, dengan $N$ jumlah getaran.
\begin{equation}\label{eq:periode}
  T \approx 2\pi\sqrt{\frac{L}{g}}.
\end{equation}
Persamaan~\ref{eq:periode} berlaku untuk sudut kecil.
\end{lstlisting}
Kurung kurawal mengelompokkan pembilang, penyebut, dan isi akar. Tanda tilde menjaga kata dan nomor rujukan pada baris yang sama.
\tugas{Kompilasi dua kali. Sisipkan persamaan bernomor sebelum rumus periode dan periksa nomor rujukannya.}
\end{frame}
\begin{frame}[fragile]{Indeks, pangkat, kurung, dan pecahan}
Tambahkan \lstinline|\usepackage{amsmath,amssymb}| di prambule.
\begin{lstlisting}
\[
  T_{10}^2,\quad \frac{a+b}{c},\quad \binom{n}{k},
  \qquad \left(\frac{T_2}{T_1}\right)^2
\]
\end{lstlisting}
\[
  T_{10}^2,\quad \frac{a+b}{c},\quad \binom{n}{k},
  \qquad \left(\frac{T_2}{T_1}\right)^2
\]
Gunakan kurung kurawal untuk indeks atau pangkat lebih dari satu karakter.
\lstinline|\left| dan \lstinline|\right| membentuk pasangan pembatas otomatis.
\tugas{Tulis akar kuadrat dari rasio panjang dua tali.}
\end{frame}
\begin{frame}[fragile]{Matriks dan vektor}
\begin{lstlisting}
\[
  A = \begin{pmatrix} 1 & 2 \\ 3 & 4 \end{pmatrix},
  \qquad
  \mathbf{t} = \begin{bmatrix} T_1 \\ T_2 \end{bmatrix}
\]
\end{lstlisting}
\[
A=\begin{pmatrix}1&2\\3&4\end{pmatrix},\qquad
\mathbf{t}=\begin{bmatrix}T_1\\T_2\end{bmatrix}
\]
\lstinline|&| memisahkan kolom dan \lstinline|\\| mengakhiri baris.
\lstinline|pmatrix| memakai tanda kurung; \lstinline|bmatrix| memakai kurung siku.
\tugas{Ubah matriks menjadi dua baris dan tiga kolom.}
\end{frame}
\begin{frame}[fragile]{Meratakan beberapa persamaan}
\begin{lstlisting}
\begin{align}
  T_1 &= \frac{12{,}7}{10} = 1{,}27 \label{eq:t1}\\
  g_1 &= \frac{4\pi^2 L_1}{T_1^2}
       \approx 9{,}79 \label{eq:g1}
\end{align}
\end{lstlisting}
\begin{align*}
T_1 &= \frac{12{,}7}{10} = 1{,}27\\
g_1 &= \frac{4\pi^2 L_1}{T_1^2}\approx9{,}79
\end{align*}
Tempatkan \lstinline|&| sebelum tanda sama dengan.
\lstinline|align| memberi nomor setiap baris; \lstinline|align*| menghilangkannya.
Jangan bungkus \lstinline|align| dengan tanda dolar atau \lstinline|\[...\]|.
\end{frame}
\begin{frame}[fragile]{Operator, jumlah, integral, dan limit}
\begin{lstlisting}
\[
  \bar{T}=\frac{1}{n}\sum_{i=1}^{n}T_i,\qquad
  \int_0^t v(\tau)\,\mathrm{d}\tau,\qquad
  \lim_{\theta\to0}\frac{\sin\theta}{\theta}=1
\]
\end{lstlisting}
\[
\bar{T}=\frac{1}{n}\sum_{i=1}^{n}T_i,\qquad
\int_0^t v(\tau)\,\mathrm{d}\tau,\qquad
\lim_{\theta\to0}\frac{\sin\theta}{\theta}=1
\]
Gunakan \lstinline|\sin|, \lstinline|\log|, dan \lstinline|\lim|
agar nama operator tegak dengan spasi yang sesuai.
\tugas{Tulis rata-rata tiga pengukuran periode dengan notasi jumlah.}
\end{frame}
\begin{frame}[fragile]{Huruf, simbol, dan spasi matematika}
\begin{lstlisting}
$\alpha,\ \theta,\ \Delta,\ \pi$
$\mathbf{x}\in\mathbb{R}^n$
$T = 1{,}27\,\mathrm{s}$
\end{lstlisting}
$\alpha,\ \theta,\ \Delta,\ \pi$ \qquad
$\mathbf{x}\in\mathbb{R}^n$ \qquad $T=1{,}27\,\mathrm{s}$
\medskip

Variabel dicetak miring; satuan menggunakan huruf tegak.
\lstinline|\,| memberi spasi tipis. Kurung pada \lstinline|{,}|
menghindari spasi tanda baca pada desimal.
Rumus terpisah memberi lebih banyak ruang daripada rumus dalam kalimat.
\end{frame}
\begin{frame}[fragile]{Dari sepuluh getaran ke satu periode}
\centering
\begin{tabular}{rrrr}\toprule
$L$ (m)&$N$&$t$ (s)&$T=t/N$ (s)\\\midrule
0,40&10&12,7&1,27\\
0,60&10&15,6&1,56\\
0,80&10&18,0&1,80\\\bottomrule
\end{tabular}
\par\medskip\raggedright
Data ilustratif untuk latihan. Dua baris mana yang dapat digunakan untuk menguji prediksi saat panjang tali dilipatduakan?
\tugas{Bandingkan $1{,}80/1{,}27$ dengan $\sqrt{2}$. Simpan hasilnya untuk pembahasan.}
\end{frame}
\begin{frame}[fragile]{Struktur tabel dan nomor rujukannya}
\begin{lstlisting}
\begin{table}[htbp]
  \centering
  \caption{Periode untuk setiap panjang tali}\label{tab:data}
  \begin{tabular}{rrrr}
    $L$ (m) & $N$ & $t$ (s) & $T$ (s) \\ \hline
    0,40 & 10 & 12,7 & 1,27 \\
  \end{tabular}
\end{table}
Data pada Tabel~\ref{tab:data} menunjukkan ...
\end{lstlisting}
\tugas{Lengkapi dua baris berikutnya. Setiap baris data memiliki empat sel yang dipisahkan oleh tiga tanda \texttt{\&}.}
\end{frame}
\begin{frame}[fragile]{Gambar alat yang dapat disunting}
Muat \lstinline|\usepackage{tikz}| di prambule. Tempatkan kode berikut pada bagian metode.
\begin{lstlisting}
\begin{figure}[htbp]
  \centering
  \begin{tikzpicture}
    \draw (-1,0)--(1,0);
    \draw[dashed] (0,0)--(0,-2);
    \draw (0,0)--(-65:2) node[midway,right] {$L$};
    \fill (-65:2) circle (0.1);
  \end{tikzpicture}
  \caption{Pengukuran panjang ayunan}\label{fig:alat}
\end{figure}
\end{lstlisting}
\end{frame}
\begin{frame}[fragile]{Memasukkan gambar dari berkas}
Muat \lstinline|\usepackage{graphicx}| di prambule.
Unggah gambar sendiri dan gunakan nama berkas yang sesuai.
\begin{lstlisting}
\begin{figure}[htbp]
  \centering
  \includegraphics[width=.6\linewidth]{figures/alat.png}
  \caption{Perangkat pengukuran}\label{fig:foto}
\end{figure}
Lihat Gambar~\ref{fig:foto}.
\end{lstlisting}
PDF cocok untuk gambar vektor; PNG atau JPEG untuk gambar raster.
Sesuaikan ukuran dengan \lstinline|\linewidth| agar mengikuti lebar teks.
\end{frame}
\begin{frame}[fragile]{Bagaimana LaTeX menempatkan gambar?}
Lingkungan \lstinline|figure| dan \lstinline|table| memberi LaTeX ruang untuk memilih letak. Opsi \lstinline|[htbp]| mengizinkan posisi di sekitar teks, atas halaman, bawah halaman, atau halaman khusus objek.

\medskip
Letakkan \lstinline|\label| sesudah \lstinline|\caption| agar kuncinya merekam nomor gambar atau tabel.
\tugas{Tulis kalimat yang menjelaskan pengukuran $L$ dan merujuk \texttt{fig:alat}. Tambahkan satu paragraf di depan gambar, lalu periksa rujukannya.}
\end{frame}
\begin{frame}[fragile]{Sitasi pertama dan daftar pustaka}
\begin{lstlisting}
Penomoran otomatis dijelaskan dalam manual \cite{latexref}.

\begin{thebibliography}{9}
  \bibitem{latexref}
  \emph{LaTeX2e: An unofficial reference manual}.
  Juni 2026, Bab 7.
\end{thebibliography}
\end{lstlisting}
Ketik sitasi di bagian yang membahas penomoran. Tempatkan daftar pustaka sebelum \lstinline|\end{document}|. Kompilasi kembali hingga nomor sitasi muncul.
\end{frame}
\section{Sesi 3. Proyek dokumen akademik lengkap}
\begin{frame}{Proyek dokumen akademik lengkap}
\relax{\large Sesi 3 \quad 13.00--14.15 \quad 75 menit}\par\bigskip
15 menit menyiapkan struktur dan metode\\[3pt]
25 menit menggabungkan persamaan, tabel, dan gambar\\[3pt]
20 menit menulis pembahasan, sitasi, dan daftar pustaka\\[3pt]
15 menit menelaah dan merevisi laporan
\medskip

\textbf{Hasil sesi}\quad Laporan pendek dengan teks, bagian, persamaan,
tabel, gambar, sitasi, dan daftar pustaka.
\end{frame}
\begin{frame}[fragile]{Menit 0--15 / Kerangka laporan}
Lanjutkan berkas pagi ini. Gunakan kelas \lstinline|article| dan
paket yang memang dibutuhkan. Setelah \lstinline|\maketitle|, tulis:
\begin{lstlisting}
\begin{abstract}
Kami membandingkan periode ayunan pada beberapa panjang tali.
\end{abstract}
\section{Pendahuluan}\label{sec:pendahuluan}
\section{Metode}\label{sec:metode}
\section{Hasil dan pembahasan}\label{sec:hasil}
\section{Kesimpulan}\label{sec:kesimpulan}
\end{lstlisting}
\tugas{Nyatakan pertanyaan, asumsi sudut kecil, dan cara mengukur satu periode.}
\end{frame}
\begin{frame}{Menit 15--40 / Data, model, dan gambar}
\begin{enumerate}
\item Masukkan model periode, definisi variabel, dan satuannya.
\item Lengkapi tabel tiga panjang tali dengan periode dan estimasi
$g=4\pi^2L/T^2$.
\item Masukkan diagram atau foto alat; beri keterangan dan label.
\item Rujuk persamaan, tabel, dan gambar dari paragraf pembahasan.
\end{enumerate}
\medskip
\textbf{Titik periksa}\quad Untuk $L=0{,}60$ m dan $T=1{,}56$ s,
estimasi $g$ sekitar $9{,}73$ m/s$^2$.
\tugas{Nyatakan bahwa data latihan bersifat ilustratif.}
\end{frame}
\begin{frame}[fragile]{Contoh perhitungan yang dapat diperiksa}
\begin{lstlisting}
\begin{align*}
  T_1 &= \frac{12{,}7}{10}\,\mathrm{s}
       = 1{,}27\,\mathrm{s},\\
  g_1 &= \frac{4\pi^2 L_1}{T_1^2}
       \approx 9{,}79\,\mathrm{m/s^2}.
\end{align*}
Data pada Tabel~\ref{tab:data} mendukung ...
\end{lstlisting}
Jelaskan arti hasil perhitungan, bukan hanya menampilkan rumus.
\tugas{Bandingkan $1{,}80/1{,}27$ dengan $\sqrt{2}$.
Apakah periode menjadi dua kali lipat ketika panjang dilipatduakan?}
\end{frame}
\begin{frame}{Menit 40--60 / Pembahasan dan sumber}
\begin{itemize}
\item Tulis satu paragraf yang membandingkan prediksi model dan data.
\item Jelaskan satu keterbatasan: waktu reaksi, pengukuran panjang,
atau simpangan yang terlalu besar.
\item Tulis kesimpulan yang menjawab pertanyaan penelitian.
\item Sitasi sumber yang benar-benar digunakan dan lengkapi daftar pustaka.
\end{itemize}
\medskip
Gunakan kunci rujukan, bukan nomor yang diketik manual.
\tugas{Baca setiap sitasi dan pastikan entri sumbernya sesuai.}
\end{frame}
\begin{frame}{Menit 60--75 / Telaah berpasangan}
\begin{enumerate}
\item \textbf{5 menit:} tukar sumber beserta gambar, lalu kompilasi.
\item \textbf{5 menit:} periksa satu perhitungan dan semua rujukan;
beri satu koreksi isi serta satu koreksi tampilan.
\item \textbf{5 menit:} terapkan koreksi dan simpan versi laporan.
\end{enumerate}
\medskip
\textbf{Kriteria selesai}\quad Ketujuh elemen hadir, rujukan terisi,
sumber dapat dikompilasi, dan kesimpulan menjawab pertanyaan.
\medskip

Laporan ini menjadi bahan penataan ulang pada sesi berikutnya.
\end{frame}
\section{Sesi 4. Manajemen proyek LaTeX untuk dokumen akademik besar}
\begin{frame}{Manajemen proyek dokumen akademik besar}
\relax{\large Sesi 4 \quad 14.30--16.00 \quad 90 menit}\par\bigskip
25 menit mengatur folder, berkas utama, dan berkas bagian\\[3pt]
20 menit memusatkan prambule dan definisi yang dipakai ulang\\[3pt]
20 menit mengelola bibliografi dan rujukan antarberkas\\[3pt]
15 menit memeriksa detail dan mendiagnosis masalah\\[3pt]
10 menit menelaah proyek dan menyerahkan hasil
\medskip

\textbf{Hasil sesi}\quad Proyek beberapa berkas yang mudah dikembangkan
menjadi laporan panjang, skripsi, atau buku.
\end{frame}
\begin{frame}{Kapan dokumen perlu dipisah?}
\begin{itemize}
\item Naskah memiliki banyak bagian, gambar, dan referensi.
\item Beberapa orang mengerjakan bagian berbeda.
\item Pengaturan tampilan yang sama muncul berulang kali.
\item Menemukan sumber kesalahan mulai memerlukan waktu lama.
\end{itemize}
\medskip
Pertahankan satu berkas utama. Pisahkan isi menurut tanggung jawab:
pengaturan, bagian naskah, gambar, dan data bibliografi.
\tugas{Duplikasi proyek sesi 3 sebagai titik pemulihan sebelum menata ulang.}
\end{frame}
\begin{frame}[fragile]{Struktur folder yang sederhana}
\begin{lstlisting}[basicstyle=\ttfamily\small\color{codeink}]
main.tex
config/
  preamble.tex
sections/
  pendahuluan.tex
  metode.tex
  hasil.tex
  kesimpulan.tex
figures/
  alat.png
references.bib
\end{lstlisting}
Gunakan nama singkat, konsisten, dan tanpa spasi.
Berkas gambar harus benar-benar tersedia pada jalur yang dipanggil.
\end{frame}
\begin{frame}[fragile]{Berkas utama sebagai pengendali}
\begin{lstlisting}
\documentclass[11pt,a4paper]{article}
\input{config/preamble}
\title{Mengukur periode ayunan}
\author{Nama peserta}
\begin{document}
\maketitle
\tableofcontents
\input{sections/pendahuluan}
\input{sections/metode}
\input{sections/hasil}
\input{sections/kesimpulan}
\printbibliography
\end{document}
\end{lstlisting}
\end{frame}
\begin{frame}[fragile]{Isi berkas bagian}
\small
Contoh isi \texttt{sections/metode.tex}:
\begin{lstlisting}
\section{Metode}\label{sec:metode}
Panjang tali diukur dari titik gantung ke pusat beban.
Waktu untuk sepuluh getaran dibagi sepuluh.

Hasilnya disajikan pada Bagian~\ref{sec:hasil}.
\end{lstlisting}
Berkas bagian tidak memuat \lstinline|\documentclass|,
\lstinline|\begin{document}|, atau \lstinline|\end{document}|.
\medskip

Kompilasi \texttt{main.tex}. Semua label di bagian yang dimuat
berbagi ruang nama yang sama; setiap label harus unik.
\tugas{Pindahkan empat bagian laporan dan bandingkan PDF sebelum serta sesudahnya.}
\end{frame}
\begin{frame}[fragile]{Memilih input atau include}
\small
\begin{description}
\item[\texttt{input}] Menyisipkan isi berkas pada tempat pemanggilan,
tanpa memulai halaman baru. Cocok untuk bagian artikel dan konfigurasi.
\item[\texttt{include}] Memisahkan bagian dengan pergantian halaman
dan berkas bantu tersendiri. Cocok untuk bab laporan atau buku.
\end{description}
\begin{lstlisting}
% Kelas report/book; kompilasi semua bab terlebih dahulu.
\includeonly{chapters/metode}
% Di dalam document:
\include{chapters/pendahuluan}
\include{chapters/metode}
\end{lstlisting}
\lstinline|\includeonly| memerlukan berkas bantu kompilasi lengkap
untuk mempertahankan rujukan bab yang dilewati.
\end{frame}
\begin{frame}[fragile]{Satu tempat untuk pengaturan bersama}
Isi \texttt{config/preamble.tex}:
\begin{lstlisting}
\usepackage[indonesian]{babel}
\usepackage{amsmath,amssymb,graphicx,booktabs}
\usepackage{tikz} % jika diagram TikZ digunakan
\usepackage{standalone}
\usepackage[backend=biber,style=numeric]{biblatex}
\addbibresource{references.bib}
\usepackage{hyperref}
\graphicspath{{figures/}}
\newcommand{\periode}[1]{T_{#1}}
\newcommand{\satuan}[1]{\,\mathrm{#1}}
\end{lstlisting}
Muat paket sekali dengan opsi yang konsisten. Pada contoh ini,
\lstinline|hyperref| dimuat setelah paket lainnya.
Jalur berkas dihitung dari direktori kompilasi berkas utama.
\end{frame}
\begin{frame}[fragile]{Perintah ulang pakai untuk detail konsisten}
Definisi di prambule:
\begin{lstlisting}
\newcommand{\periode}[1]{T_{#1}}
\newcommand{\satuan}[1]{\,\mathrm{#1}}
\end{lstlisting}
Penggunaan di berkas bagian:
\begin{lstlisting}
$\periode{1}=1{,}27\satuan{s}$
$\periode{2}=1{,}56\satuan{s}$
\end{lstlisting}
Argumen \lstinline|#1| diganti dengan isi kurung.
Kedua perintah contoh dipakai di dalam mode matematika.
\tugas{Ganti dua penulisan periode; ubah definisi sekali dan periksa keduanya.}
\end{frame}
\begin{frame}[fragile]{Lingkungan untuk catatan yang berulang}
Definisikan sekali di prambule:
\begin{lstlisting}
\newenvironment{catatan}
  {\begin{quote}\small\textbf{Catatan.}\ }
  {\end{quote}}
\end{lstlisting}
Gunakan di bagian mana pun:
\begin{lstlisting}
\begin{catatan}
Data pada tabel digunakan untuk latihan penulisan.
\end{catatan}
\end{lstlisting}
\tugas{Tambahkan dua catatan dan ubah gaya bersama dari satu definisi.}
\end{frame}
\begin{frame}[fragile]{Memindahkan referensi ke basis data}
\small
Buat \texttt{references.bib}. Contoh berikut adalah entri latihan
fiktif; ganti dengan sumber yang benar-benar digunakan.
\begin{lstlisting}
@book{contoh2026,
  author    = {Penulis, Nama},
  title     = {Contoh Buku untuk Latihan},
  year      = {2026},
  publisher = {Penerbit Contoh}
}
\end{lstlisting}
Di berkas isi, gunakan \lstinline|\cite{contoh2026}|.
Di akhir berkas utama, gunakan \lstinline|\printbibliography|.
\medskip

Hapus lingkungan \lstinline|thebibliography| lama setelah entri sumber
dipindahkan. Kunci setiap entri harus unik.
\end{frame}
\begin{frame}[fragile]{Alur kompilasi bibliografi}
\small
Contoh sesi ini menggunakan \textbf{biblatex + Biber}.
Dengan pengaturan \lstinline|backend=biber|, urutan manual lokal adalah:
\begin{lstlisting}[style=terminal]
pdflatex main
biber main
pdflatex main
pdflatex main
\end{lstlisting}
\lstinline|latexmk -pdf main.tex| biasanya mengelola urutan ini.
Di Overleaf, kompilasi proyek dari \texttt{main.tex}.
\medskip

Template lama dapat menggunakan BibTeX dengan
\lstinline|\bibliographystyle| dan \lstinline|\bibliography|.
Ikuti sistem template; jangan gabungkan dua sistem dalam proyek yang sama.
\end{frame}
\begin{frame}[fragile]{Rujukan antarberkas dan penamaan label}
\begin{lstlisting}
% sections/metode.tex
\section{Metode}\label{sec:metode}
% sections/hasil.tex
Metode pada Bagian~\ref{sec:metode} digunakan untuk ...
Persamaan~\eqref{eq:periode} memprediksi ...
Data terdapat pada Tabel~\ref{tab:data}.
\end{lstlisting}
Gunakan awalan yang membantu pencarian:
\texttt{sec:}, \texttt{eq:}, \texttt{fig:}, dan \texttt{tab:}.
Label gambar dan tabel ditempatkan setelah keterangan.
\tugas{Sisipkan satu persamaan baru dan pastikan rujukan nomor tetap benar.}
\end{frame}
\begin{frame}[fragile]{Detail kecil: nomor, daftar, dan spasi}
\begin{lstlisting}
\tableofcontents
\listoffigures
\listoftables
% Di bagian yang sesuai:
Lihat Gambar~\ref{fig:alat}.
$g=9{,}81\,\mathrm{m/s^2}$
\end{lstlisting}
\begin{itemize}
\item Biarkan judul bagian dan keterangan mengisi daftar otomatis.
\item Seragamkan desimal, satuan, singkatan, dan tanda baca.
\item Jangan mengatur posisi dengan deretan spasi atau pemutus baris.
\item Kompilasi ulang setelah perubahan judul atau penomoran.
\end{itemize}
\end{frame}
\begin{frame}{Template dan kelas dokumen}
\begin{itemize}
\item Mulai dari contoh minimal yang disertakan dalam template.
\item Kenali berkas utama, kelas \texttt{.cls}, paket \texttt{.sty},
mesin kompilasi, dan sistem bibliografinya.
\item Pindahkan isi secara bertahap dan kompilasi setiap tahap.
\item Simpan penyesuaian di prambule atau berkas konfigurasi sendiri.
\end{itemize}
\medskip
Jangan mengubah \lstinline|article| menjadi \lstinline|book| tanpa
meninjau hierarki bab, halaman judul, dan penomoran.
Pengaturan tata letak mendalam dilanjutkan pada Hari 3.
\end{frame}
\begin{frame}[fragile]{Mendiagnosis proyek besar}
\small
\begin{tabular}{@{}p{.36\linewidth}p{.59\linewidth}@{}}
\toprule
Gejala & Pemeriksaan pertama\\\midrule
Berkas tidak ditemukan & Cocokkan nama, huruf besar/kecil, dan jalur dari berkas utama.\\
Perintah tidak dikenal & Periksa ejaan, paket, dan urutan definisi.\\
Label ganda atau \texttt{??} & Cari label yang berulang atau salah, lalu kompilasi ulang.\\
Sitasi tidak terisi & Periksa kunci, berkas \texttt{.bib}, dan log Biber.\\
\texttt{Overfull} & Periksa baris, tabel, atau gambar yang terlalu lebar.\\
\bottomrule
\end{tabular}
\tugas{Pada salinan proyek, ubah satu jalur berkas. Baca kesalahan pertama,
temukan penyebabnya, lalu pulihkan jalur tersebut.}
\end{frame}
\begin{frame}{Menit 80--90 / Serah-terima proyek}
\begin{enumerate}
\item \textbf{4 menit:} rekan membuka salinan proyek lengkap dan
mengompilasi berkas utama.
\item \textbf{3 menit:} periksa daftar isi, nomor, rujukan, sitasi,
keterangan gambar, serta teks yang melewati margin.
\item \textbf{3 menit:} terapkan koreksi dan simpan sumber beserta PDF.
\end{enumerate}
\medskip
Sertakan seluruh berkas bagian, konfigurasi, gambar, dan
\texttt{.bib}. Catat berkas utama dan mesin kompilasi yang digunakan.
\medskip

Simpan tabel dan kesimpulan untuk presentasi Beamer pada Hari 2.
\end{frame}
\begin{frame}{Capaian akhir Hari 1}
\begin{itemize}
\item Membuat dan memperbaiki dokumen LaTeX dasar.
\item Menulis elemen akademik dengan nomor dan rujukan otomatis.
\item Menggabungkan elemen tersebut menjadi laporan utuh.
\item Memisahkan proyek menjadi berkas yang terorganisasi.
\item Menjaga konsistensi detail dari definisi dan pengaturan bersama.
\end{itemize}
\medskip
\textbf{Dua keluaran}\quad PDF yang siap dibaca dan proyek sumber
yang dapat dikompilasi ulang oleh orang lain.
\end{frame}
\end{document}
